
The numerical calculation of bounce coefficients requires an integration
over $\theta $ which can be expressed symbolically as
\begin{equation}
I\left( \psi ,\xi _{0}\right) =\int_{\theta _{\min }}^{\theta _{\max }}\frac{%
d\theta }{2\pi }F\left( \psi ,\theta ,\xi _{0};B_{R},B_{Z},B_{\phi
},R,Z\right)
\end{equation}
where $B_{R},B_{Z},B_{\phi },R,Z$ are functions of $\left( \psi ,\theta
\right) $. They are given on a uniform grid of $N_{\theta }$ points in $%
\theta $
\begin{equation}
\theta _{j}=\frac{2\pi j}{N_{\theta }-1},\;j=0,1,\cdots ,N_{\theta }-1
\end{equation}

\paragraph{Domain of Integration}

It is very important to account for the entire bounce path of the particle,
including in particular the tip of banana orbits near $\theta _{T\min }$ and
$\theta _{T\max }$. The contribution of these banana tips is often larger
than the $d\theta =2\pi /(N_{\theta }-1)$ accuracy level, because $F\left(
\theta \right) $ can become very large near the turning points. This is true
for example in the calculation of $\lambda $, since $F\left( \theta \right)
\sim $ $1/\xi $ and $\xi \rightarrow 0$ at the turning points. It is
therefore crucial to perform the integration up to $\theta _{T\min }$ and $%
\theta _{T\max }$. However, these boundaries are defined by (\ref{thetamin})
and (\ref{thetamax})
\begin{equation}
B\left( \psi ,\theta _{T}\right) =B_{b}\left( \psi ,\xi _{0}\right) =\frac{%
B_{0}\left( \psi \right) }{1-\xi _{0}^{2}}
\end{equation}
which in general do not coincide with a grid points. In order to calculate $%
\theta _{T}$, we impose that the values of the data $B_{R},B_{Z},B_{\phi
},R,Z$ in $\theta _{T}$ be obtained by linear interpolation, while the value
of $B\left( \psi ,\theta _{T}\right) $ is obtained from (\ref{Btot}).

We consider the magnetic field $B_{b}\left( \psi ,\xi _{0}\right) $ at the
turning point $\theta _{T\min }$ to be located between the two (consecutive)
values $B_{1}\left( \psi ,\theta _{1}\right) $ and $B_{2}\left( \psi ,\theta
_{2}\right) $ on the $\left( \psi ,\theta \right) $ grid. These values are
determined from the three components issued by the equilibrium code simply
by
\begin{eqnarray}
B_{1}\left( \psi ,\theta _{1}\right)  &=&\sqrt{B_{R}^{2}\left( \psi ,\theta
_{1}\right) +B_{Z}^{2}\left( \psi ,\theta _{1}\right) +B_{\phi }^{2}\left(
\psi ,\theta _{1}\right) } \\
B_{2}\left( \psi ,\theta _{2}\right)  &=&\sqrt{B_{R}^{2}\left( \psi ,\theta
_{2}\right) +B_{Z}^{2}\left( \psi ,\theta _{2}\right) +B_{\phi }^{2}\left(
\psi ,\theta _{2}\right) }
\end{eqnarray}

We choose to define the values of $B_{R}$, $B_{Z}\ $and $B_{\phi }$ at the
location $\theta _{T}$ by linear interpolation:
\begin{equation}
B_{i}\left( \psi ,\theta _{T}\right) =B_{i1}+\frac{\left( \theta _{T}-\theta
_{1}\right) }{\left( \theta _{2}-\theta _{1}\right) }\left(
B_{i2}-B_{i1}\right)
\end{equation}
where $i=R,Z,\phi $. Then, the location $\theta _{T}$ of the turning point
can be calculated by requiring that the relation
\begin{equation}
B_{b}\left( \psi ,\xi _{0}\right) =B\left( \psi ,\theta _{T}\right) =\sqrt{%
B_{R}^{2}\left( \psi ,\theta _{T}\right) +B_{Z}^{2}\left( \psi ,\theta
_{T}\right) +B_{\phi }^{2}\left( \psi ,\theta _{T}\right) }
\end{equation}
be satisfied. This implies
\begin{equation}
B_{R}^{2}\left( \psi ,\theta _{T}\right) +B_{Z}^{2}\left( \psi ,\theta
_{T}\right) +B_{\phi }^{2}\left( \psi ,\theta _{T}\right) =B_{b}^{2}\left(
\psi ,\xi _{0}\right)
\end{equation}
and then
\begin{equation}
\sum_{i=R,Z,\phi }\left[ B_{i1}+\frac{\left( \theta _{T}-\theta _{1}\right)
}{\left( \theta _{2}-\theta _{1}\right) }\left( B_{i2}-B_{i1}\right) \right]
^{2}-B_{b}^{2}\left( \psi ,\xi _{0}\right) =0
\end{equation}

Defining
\begin{equation}
\alpha =\frac{\left( \theta _{T}-\theta _{1}\right) }{\left( \theta
_{2}-\theta _{1}\right) }
\end{equation}
we find
\begin{equation}
\left[ \sum_{i=R,Z,\phi }\left( B_{i2}-B_{i1}\right) ^{2}\right] \alpha
^{2}+2\left[ \sum_{i=R,Z,\phi }B_{i1}\left( B_{i2}-B_{i1}\right) \right]
\alpha +\sum_{i=R,Z,\phi }B_{i1}^{2}-B_{b}^{2}\left( \psi ,\xi _{0}\right) =0
\end{equation}
which solves as
\begin{equation}
\alpha =\frac{\pm \sqrt{\left[ \sum_{i}B_{i1}\left( B_{i2}-B_{i1}\right)
\right] ^{2}-\left[ \sum_{i}\left( B_{i2}-B_{i1}\right) ^{2}\right] \left[
\sum_{i}B_{i1}^{2}-B_{b}^{2}\right] }%
-\sum_{i}B_{i1}\left( B_{i2}-B_{i1}\right) }{\sum_{i}\left(
B_{i2}-B_{i1}\right) ^{2}}
\end{equation}

We have
\begin{eqnarray}
&&\sqrt{\left[ \sum_{i}B_{i1}\left( B_{i2}-B_{i1}\right) \right] ^{2}-\left[
\sum_{i}\left( B_{i2}-B_{i1}\right) ^{2}\right] \left[
\sum_{i}B_{i1}^{2}-B_{b}^{2}\left( \psi ,\xi _{0}\right) \right] }  \nonumber
\\
&=&\sqrt{%
\begin{array}{c}
\left( B_{R1}B_{R2}+B_{Z1}B_{Z2}+B_{\phi 1}B_{\phi 2}\right)
^{2}-B_{1}^{2}B_{2}^{2} \\
+\left[ B_{2}^{2}+B_{1}^{2}-2\left( B_{R2}B_{R1}+B_{Z2}B_{Z1}+B_{\phi
2}B_{\phi 1}\right) \right] B_{b}^{2}\left( \psi ,\xi _{0}\right)
\end{array}
}  \nonumber \\
&=&\sqrt{\left( Y-B_{b}^{2}\right) ^{2}-\left( B_{b}^{2}-B_{1}^{2}\right)
\left( B_{b}^{2}-B_{2}^{2}\right) }
\end{eqnarray}
with
\begin{equation}
Y=B_{R1}B_{R2}+B_{Z1}B_{Z2}+B_{\phi 1}B_{\phi 2}
\end{equation}
so that
\begin{equation}
\alpha _{b}=\frac{\pm \sqrt{\left( Y-B_{b}^{2}\right) ^{2}-\left(
B_{b}^{2}-B_{1}^{2}\right) \left( B_{b}^{2}-B_{2}^{2}\right) }+B_{1}^{2}-Y}{%
B_{1}^{2}+B_{2}^{2}-2Y}
\end{equation}
and finally
\begin{equation}
\theta _{T}=\theta _{1}+\left( \theta _{2}-\theta _{1}\right) \frac{\pm
\sqrt{\left( Y-B_{b}^{2}\right) ^{2}-\left( B_{b}^{2}-B_{1}^{2}\right)
\left( B_{b}^{2}-B_{2}^{2}\right) }+B_{1}^{2}-Y}{B_{1}^{2}+B_{2}^{2}-2Y}
\end{equation}

We must choose (numerically) the solution that gives $0\leq \alpha _{b}\leq 1
$. Note that if the magnetic fields in points $1$ and $2$ are equal, we have
$Y=B_{1}^{2}=B_{2}^{2}=B_{b}^{2}$.

\paragraph{Numerical Integration}

Once the two turning points have been added to the $\theta $ grid, now noted
$\widetilde{\theta }_{j}$, $j=0,1,2,\cdots N_{\theta }+1$, we define the
half grid
\begin{equation}
\overline{\theta }_{k}=\frac{\left( \widetilde{\theta }_{k+1}+\widetilde{%
\theta }_{k}\right) }{2},\;k=0,1,2,\cdots N_{\theta }
\end{equation}
and calculate the discrete function
\begin{equation}
F_{k}=F\left( \psi ,\overline{\theta }_{k},\xi _{0};B_{Rk},B_{Zk},B_{\phi
k},R_{k},Z_{k}\right)
\end{equation}
where $B_{Rk},B_{Zk},B_{\phi k},R_{k},Z_{k}$ have been calculated on the
grid $\overline{\theta }_{k}$ by linear interpolation, and the step $%
\overline{d\theta }_{k}$, defined by
\begin{equation}
\overline{d\theta }_{k}=\widetilde{\theta }_{k+1}-\widetilde{\theta }
_{k},\;k=0,1,2,\cdots N_{\theta }
\end{equation}
so that the integral becomes
\begin{equation}
I\left( \psi ,\xi _{0}\right) =\frac{1}{2\pi }\sum_{k=0}^{N_{\theta }}%
\overline{d\theta }_{k}F_{k}
\end{equation}
